\chapter{Experimental Validation and System Outlook}

\section{Proposal-Stage Experimental Direction}

The simulation study establishes the design rules for injection-locking-based distributed coherent processing. The proposal-stage experimental work provides the next step: validating these ideas using physical semiconductor lasers, optical modulators, optical spectrum measurements, and coherent interference measurements. The experiments are organized around three questions. First, can a local laser be reliably locked to a remote source carrier? Second, how does modulation on the source influence the locking process? Third, can the locked laser be used in a coherent interference system for optical-domain computation?

\section{Locking-Range Measurement}

The unmodulated-source experiment measures the baseline locking range of the local injected laser. A tunable source laser is attenuated to set the injection ratio and injected into the target laser through an optical circulator. The output is monitored in both the wavelength domain and the frequency domain. This measurement establishes where the target mode dominates and how the locking boundary shifts with injection ratio.

\begin{figure}[t]
    \centering
    \includegraphics[width=0.84\textwidth]{nomodulation_setup.png}
    \caption{Experimental setup for measuring injection-locking range with an unmodulated source laser. The electro-optic modulator is retained in the path so that the optical layout remains consistent with the modulated-source experiment.}
    \label{fig:exp_nomod_setup}
\end{figure}

\begin{figure}[t]
    \centering
    \includegraphics[width=0.72\textwidth]{LR_IR.png}
    \caption{Proposal-stage locking-range measurement as a function of injection ratio. The data show that the locking range broadens as injection ratio increases and that the locking boundary is asymmetric with detuning.}
    \label{fig:exp_lr_ir}
\end{figure}

\section{Modulated Source Injection}

The modulated-source experiment introduces sidebands on the remote signal before injection. This is closer to the intended distributed computing use case, where the remote carrier carries data. As the carrier and its modulation sidebands sweep across the locking region, different spectral components can interact with the injected laser. The measurement therefore probes whether sidebands participate in locking or create unwanted nonlinear response.

\begin{figure}[t]
    \centering
    \includegraphics[width=0.84\textwidth]{modulation_setup.png}
    \caption{Experimental setup for injection locking with a modulated source laser. Modulation sidebands around the carrier can influence the locking process when they enter or approach the locking range.}
    \label{fig:exp_mod_setup}
\end{figure}

\section{Coherent Interference Experiment}

After carrier locking is established, the local laser can be remodulated and interfered with the retained remote branch. A balanced detection architecture extracts the interference cross term, while phase stabilization is used to hold the interferometer at a useful operating point. This experiment connects the laser-level locking measurements to the coherent-computing mechanism studied in simulation.

\begin{figure}[t]
    \centering
    \includegraphics[width=0.9\textwidth]{interference_sch.png}
    \caption{Experimental coherent-interference setup using an injection-locked local laser and a retained branch of the remote source signal.}
    \label{fig:exp_interference}
\end{figure}

\begin{figure}[t]
    \centering
    \includegraphics[width=0.76\textwidth]{bpd_sch.png}
    \caption{Balanced-detection concept for extracting the interference cross term used in coherent optical computation.}
    \label{fig:exp_bpd}
\end{figure}

\section{Outlook}

The experimental direction complements the simulation results by testing the same design trade-offs in a physical optical system. The eventual goal is to identify operating conditions where the local laser remains locked to the remote carrier, modulation leakage is low, and coherent interference is stable enough for analog optical computation. Future work will extend this validation to higher symbol rates, longer fiber links, polarization drift, amplified links, wavelength-division multiplexing, and integrated photonic implementations.
